%%%%%%%%%%%%%%%%%%%%%%%%
% $Autor: MansourAhmadyPhoulady, RanjeetsingTawar, HemanthPrabhakaran$
% $Pfad: Car_Licence_Plate_Identification_with_a_JetsonNano$
% $Version: 4.0.0 $
%
% !TeX encoding = utf8
%
%%%%%%%%%%%%%%%%%%%%%%%%


\chapter{Programming}

\section{Readability}
The code's readability is enhanced by several best practices:
\begin{itemize}
	\item Sequential steps guide the user through the setup, data preparation, model training, and exporting processes.
	\item Descriptive variable names and commented sections provide clarity on the purpose of each part of the code.
	\item The mixture of Python scripting and shell commands within a Python environment exemplifies a practical approach to executing a comprehensive workflow.
\end{itemize}

\section{Structure and Modules}
The workflow leverages Python's built-in modules, TensorFlow's Object Detection API, and shell commands. Key highlights include:
\begin{itemize}
	\item Use of Python's \texttt{os} and \texttt{subprocess} modules for environment setup and executing external commands.
	\item Integration of TensorFlow scripts for model training and evaluation.
	\item Commands for converting models for different deployment platforms, demonstrating the workflow's versatility.
\end{itemize}

\section{Code Snippets}

\subsection{Environment Setup}
\begin{lstlisting}[language=Python]
	import os
	import subprocess
	
	# Example of setting up paths
	paths = {'WORKSPACE_PATH': os.path.join('Tensorflow', 'workspace')}
	for path in paths.values():
	os.makedirs(path, exist_ok=True)
\end{lstlisting}

\subsection{Model Training Command}
\begin{lstlisting}[language=Python]
	# Command to start model training
	command = "python model_main_tf2.py --model_dir=path/to/model_dir --pipeline_config_path=path/to/config --num_train_steps=10000"
	subprocess.run(command.split())
\end{lstlisting}

\subsection{Dataset Preparation}
\begin{lstlisting}[language=Python]
	# Automating dataset conversion to TFRecord
	def prepare_dataset(images_dir, annotations_dir, output_dir):
	command = f"python dataset_converter.py --images_dir={images_dir} --annotations_dir={annotations_dir} --output_dir={output_dir}"
	subprocess.run(command.split())
\end{lstlisting}





\section{Parameter Handling}
This document details the handling of key parameters within the code, focusing on their roles in configuring the environment, preparing data, and controlling the training and evaluation process.


\subsection{Model Names and URL}

\begin{lstlisting}[language=Python]
	CUSTOM_MODEL_NAME = 'my_ssd_mobnet' 
	PRETRAINED_MODEL_NAME = 'ssd_mobilenet_v2_fpnlite_320x320_coco17_tpu-8'
	PRETRAINED_MODEL_URL = 'http://download.tensorflow.org/models/object_detection/tf2/20200711/ssd_mobilenet_v2_fpnlite_320x320_coco17_tpu-8.tar.gz'
\end{lstlisting}

These parameters define the custom and pretrained model names, along with the URL for downloading the pretrained model weights.

\subsection{Paths and Files}

\begin{lstlisting}[language=Python]
	paths = {'WORKSPACE_PATH': os.path.join('Tensorflow', 'workspace'), ... }
	files = {'PIPELINE_CONFIG': os.path.join('Tensorflow', 'workspace', 'models', CUSTOM_MODEL_NAME, 'pipeline.config'), ... }
\end{lstlisting}

Organizes paths and file names used throughout the setup and training process, facilitating project structure organization.

\section{Model Training and Evaluation Parameters}

\subsection{Pipeline Configuration}

\begin{lstlisting}[language=Python]
	pipeline_config = pipeline_pb2.TrainEvalPipelineConfig()
	with tf.io.gfile.GFile(files['PIPELINE_CONFIG'], "r") as f:                                                                                                                                                                                                                     
	proto_str = f.read()                                                                                                                                                                                                                                          
	text_format.Merge(proto_str, pipeline_config)  
\end{lstlisting}

Loads and sets the model architecture, input sources, and training parameters from the `.config` file.

\section{Data Preparation Parameters}

\subsection{Label Map and TFRecord Scripts}

\begin{lstlisting}[language=Python]
	labels = [{'name':'licence', 'id':1}]
	...
	!python {files['TF_RECORD_SCRIPT']} -x {os.path.join(paths['IMAGE_PATH'], 'train')} -l {files['LABELMAP']} -o {os.path.join(paths['ANNOTATION_PATH'], 'train.record')}
\end{lstlisting}

Defines classes and automates dataset conversion to the TFRecord format for TensorFlow compatibility.

\section{Training Command Parameters}

\subsection{Model Directory, Configuration Path, and Training Steps}

\begin{lstlisting}[language=Python]
	!python model_main_tf2.py --model_dir={model_dir} --pipeline_config_path={pipeline_config_path} --num_train_steps=10000
\end{lstlisting}

Specifies the directory for saving model checkpoints, the pipeline configuration path, and the number of training steps.


Understanding these parameters and their implementation within the TensorFlow Object Detection workflow allows for a tailored approach to model training and deployment, accommodating specific project needs.


\section{Error Handling}
This document outlines the addition of error handling mechanisms. Proper error handling is crucial for diagnosing and managing issues effectively, enhancing the robustness and user experience of the workflow.

\section{For Models and Dependencies}

\subsection{Error Handling for  Models}

\begin{lstlisting}[language=Python]
	import subprocess
	
	try:
	subprocess.check_call(["wget", PRETRAINED_MODEL_URL])
	except subprocess.CalledProcessError as e:
	print(f"Failed to download the pretrained model: {e}")
\end{lstlisting}

\textbf{Explanation:} This code attempts to download a pretrained model using the `wget` command. If the download fails, it catches the `CalledProcessError`, providing a clear message about the failure.

\section{Model Training}

\subsection{Error Handling During Model Training}

\begin{lstlisting}[language=Python]
	try:
	subprocess.check_call(["python", "model_main_tf2.py", f"--model_dir={model_dir}", f"--pipeline_config_path={pipeline_config_path}", "--num_train_steps=10000"])
	except subprocess.CalledProcessError as e:
	print(f"Training failed with error: {e}")
\end{lstlisting}

\textbf{Explanation:} The code snippet ensures that errors during model training, such as misconfiguration or runtime issues, are caught and reported, allowing for easier troubleshooting.

\section{TFRecord Generation}

\subsection{Error Handling for TFRecord Generation}

\begin{lstlisting}[language=Python]
	try:
	subprocess.check_call(["python", files['TF_RECORD_SCRIPT'], "-x", os.path.join(paths['IMAGE_PATH'], 'train'), "-l", files['LABELMAP'], "-o", os.path.join(paths['ANNOTATION_PATH'], 'train.record')])
	except subprocess.CalledProcessError:
	print("Failed to generate TFRecord. Please check the input paths and data format.")
\end{lstlisting}

\textbf{Explanation:} This modification catches failures in generating TFRecords, alerting the user to check their dataset's paths and formatting, thereby preventing silent failures.

\section{General Error Handling Strategy}

\subsection{Template for subprocess Calls}

\begin{lstlisting}[language=Python]
	try:
	subprocess.check_call([command, arg1, arg2])
	except subprocess.CalledProcessError as e:
	print(f"Command failed: {e}")
	except EnvironmentError as e:
	print(f"Environment setup issue: {e}")
\end{lstlisting}

\textbf{Explanation:} A general error handling template for subprocess calls. It's designed to catch and report command execution failures and broader environment issues, facilitating easier debugging and system maintenance.


Adding error handling  enhances its reliability. By ensuring that errors are caught and handled gracefully, developers can more effectively diagnose and resolve issues, leading to a smoother development and deployment process.


\section{Message Handling}
This document details the implementation of enhanced message handling in the code. Providing users with clear, informative messages significantly improves usability and aids in troubleshooting.


\section{Environment Setup and Dependency Installation}

Enhanced messaging during the setup phase informs users about the installation status of required packages.

\subsection{Installing Required Packages}

\begin{lstlisting}[language=Python]
	def install_package(package):
	try:
	subprocess.check_call(["pip", "install", package])
	print(f"Successfully installed {package}.")
	except subprocess.CalledProcessError:
	print(f"Failed to install {package}. Please check your environment and try again.")
\end{lstlisting}

\textbf{Explanation:} This function wraps the package installation in a try-except block, catching any errors that occur and providing specific feedback to the user, enhancing the setup experience.

\section{Downloading Pretrained Models}

Informative messages during model downloading help users understand the process and troubleshoot if necessary.

\subsection{Model Download}

\begin{lstlisting}[language=Python]
	def download_model(url):
	try:
	print("Starting download...")
	subprocess.check_call(["wget", url])
	print("Download completed successfully.")
	except subprocess.CalledProcessError:
	print("Download failed. Please check the URL and your internet connection.")
\end{lstlisting}

\textbf{Explanation:} This code snippet includes messages that inform the user when the download process starts and ends, providing guidance on what to check in case of a download failure.

\section{Training the Model}

Messaging during model training offers insights into the training progress and outcomes.

\subsection{Training Process}

\begin{lstlisting}[language=Python]
	def train_model(model_dir, pipeline_config_path, num_train_steps):
	try:
	print("Training started. This may take a while...")
	subprocess.check_call(["python", "model_main_tf2.py", f"--model_dir={model_dir}", f"--pipeline_config_path={pipeline_config_path}", f"--num_train_steps={num_train_steps}"])
	print("Training completed successfully.")
	except subprocess.CalledProcessError:
	print("Training failed. Please check the configuration and dataset.")
\end{lstlisting}

\textbf{Explanation:} The addition of start and completion messages for the training process, along with error notifications, keeps the user informed and aids in troubleshooting training-related issues.



The incorporation of enhanced message handling within the workflow provides users with a more informative and engaging experience. By systematically implementing clear and actionable messaging, users can more effectively navigate the setup, training, and utilization of object detection models.

